\customchapter{RESUMEN}

La documentación digital del patrimonio construido depende de levantamientos
láser cuya interpretación sigue siendo manual. Este trabajo propone un modelo
híbrido que sustituye el descriptor geométrico del mejor método de descubrimiento
de clases novedosas por un codificador auto-supervisado, y mide la mejora sobre
el levantamiento de la Catedral de Lima.

\begin{itemize}
  \item \textbf{Introducción a la Investigación:} Un escáner láser registra dónde
  hay materia, pero no qué elemento representa cada punto. Los métodos de
  segmentación semántica tridimensional usan una lista cerrada de categorías y
  asignan todo elemento imprevisto a la clase más parecida. El objetivo es
  mejorar de forma medible el reconocimiento de los elementos sin etiqueta.

  \item \textbf{Marco Metodológico:} La investigación es aplicada, cuantitativa y
  de diseño experimental comparativo. Primero se diagnostica el mejor método de
  descubrimiento de clases novedosas reanalizando su ablación; luego se sustituye
  su componente más débil por un codificador auto-supervisado. La evaluación usa
  el levantamiento de la Catedral de Lima, en cinco secciones, contra una línea
  base de tres corridas.

  \item \textbf{Resultados:} El método base logra entre 11,8 y 47,5 puntos de
  intersección sobre unión en las clases no vistas, frente a 49,8 del supervisado.
  Su ablación da al descriptor geométrico solo 4,2 puntos sobre una base de 23,2,
  menos de la mitad de los 9,6 del hipergrafo. Eliminar esa dependencia eleva un
  descriptor congelado de 5,6 a 72,5 puntos. Se define además una medida de
  brecha cerrada, pues el sesgo hacia las clases mayoritarias cuesta hasta 15,34
  puntos.

  \item \textbf{Conclusiones:} El desequilibrio del método base lo documentan sus
  propios autores y sigue sin corregirse, pero queda por establecer si persiste
  sin etiquetas. Cada punto de mejora es corrección manual que el especialista
  deja de hacer.
\end{itemize}

\vspace{1em}

\noindent \textbf{Palabras clave:}\\
\noindent Nubes de puntos; Segmentación semántica tridimensional; Descubrimiento
de clases novedosas; Aprendizaje auto-supervisado; Hipergrafos; Patrimonio
arquitectónico
